\chapter{Recommendations and Conclusions}
\label{ch:conclusion}

\section{Summary of Contributions}

This dissertation set out to evaluate whether a free-tier, open-weight LLM (Llama 3.1, 8B parameters, via the Groq API) can generate personalised UK credit-card recommendations competitive with a conventional content-based baseline, and how prompting strategy affects the quality, factual reliability, explanation quality, and consistency of its output. Against a fixed catalogue of 20 real UK credit cards and 300 ONS-calibrated synthetic user profiles spanning eight consumer archetypes, three prompting strategies (zero-shot, structured, few-shot) were evaluated across 900 LLM inference calls, scored against a content-based cosine-similarity baseline and a four-criteria evaluation framework (recommendation quality, factual correctness, explanation quality, consistency).

The central empirical finding is that prompting strategy has a large, statistically significant effect on recommendation quality in this domain ($H = 60.70$, $p = 6.6\times10^{-14}$), but the effect runs in the opposite direction to what the general prompt-engineering literature would predict: zero-shot prompting outperformed both structured and few-shot prompting on recommendation-baseline overlap and top-1 accuracy. At the same time, structured prompting produced the most factually reliable responses (73.7\% entirely error-free, against 55.0\% for zero-shot and only 24.0\% for few-shot), and few-shot prompting --- despite generating the highest volume of specific factual claims --- was the least reliable strategy on both recommendation quality and factual correctness. Most strikingly, few-shot prompting produced the \textit{highest}-rated explanations by an independent LLM judge (mean 2.907/3, against 2.730 for zero-shot and 2.650 for structured, $p = 5.1\times10^{-13}$) despite being the least accurate and least factually reliable strategy on every other measure --- demonstrating concretely that explanation quality, judged in isolation, is not a proxy for correctness. No single strategy dominates across every criterion, which is the core methodological argument this dissertation makes: a recommendation task evaluated on a single accuracy-like metric, or on explanation quality alone, would have reached a materially incomplete, and potentially misleading, conclusion about which prompting strategy is ``best.''

\section{Answering the Research Questions}

\begin{description}
    \item[RQ1] LLM-generated recommendations showed substantial but partial overlap with the content-based baseline (mean overlap 0.517--0.687 depending on strategy), confirming an LLM can approximate a purpose-built recommender without task-specific training, while not fully replicating its output.
    \item[RQ2] Prompting strategy produces a statistically significant difference in recommendation quality (Kruskal-Wallis $p < 0.0001$), with zero-shot significantly outperforming both other strategies (paired Wilcoxon $p < 0.0001$ in both comparisons) and structured vs.\ few-shot not reaching significance ($p = 0.0588$).
    \item[RQ3] Factual reliability of stated reasoning varies substantially and independently of recommendation quality: structured prompting is most reliable (73.7\% error-free), few-shot least (24.0\% error-free). Explanation quality (LLM-judged, 1--3 scale) runs in the opposite direction to factual reliability: few-shot scores highest (2.907/3), structured lowest (2.650/3), both differences statistically significant. Recommendation accuracy, factual trustworthiness, and perceived explanation quality are therefore three separate properties of a response, not three views of the same thing, and must be evaluated separately.
    \item[RQ4] Recommendation stability across repeated runs is numerically highest for zero-shot (mean pairwise Jaccard 0.826) and lowest for structured (0.776), with few-shot in between (0.808), though this difference is not statistically significant on the stratified sample tested (Kruskal-Wallis $p = 0.68$). Zero-shot's numerical lead on both accuracy (RQ2) and stability (RQ4) is internally consistent; structured's combination of high factual reliability (RQ3) but the lowest numerical stability (RQ4) is a genuine tension worth further investigation at greater sample size.
\end{description}

\section{Recommendations}

For practitioners considering an LLM-based recommender for a similarly bounded, factually-checkable product-recommendation task: this study's results suggest that (i) prompting strategy should be treated as an empirical question to test against the specific task and model, not assumed from general-purpose benchmarks, since the direction of effect found here (zero-shot best) is the reverse of the common expectation; (ii) recommendation-quality, factual-reliability, and explanation-quality metrics should be evaluated and reported separately rather than assuming any one is a proxy for the others, since a strategy that looks best on one can be worst on another, as shown starkly by few-shot prompting scoring best on explanation quality while scoring worst on both accuracy and factual reliability in this study; (iii) if factual reliability of the stated reasoning matters for a given application (e.g.\ a live consumer-facing recommender, where a factually wrong claim about a fee or reward rate is a real-world liability), structured output prompting is worth the trade-off in raw recommendation-overlap performance found here; and (iv) explanation quality alone --- whether judged by an LLM or a human reader without access to ground truth --- should not be used as a stand-in for factual accuracy when a deterministic check against real product data is possible, since this study found the two to move in opposite directions.

For future work building directly on this dissertation: the model-scale hypothesis proposed in Section~5.1.2 (that limited attention capacity at 8B parameters specifically disadvantages longer structured/few-shot prompts) could be tested directly by repeating this exact experimental design with a larger model (e.g.\ a 70B-parameter model, also available on Groq's free tier) to see whether the zero-shot advantage persists, narrows, or reverses at scale. The consistency evaluation could be extended from the stratified 24-profile subsample used here to the full 300-profile set, infrastructure permitting. Finally, the deterministic factual-correctness checker built for this study (\texttt{src/factual\_correctness.py}) is itself a reusable contribution: it could be extended to a wider card catalogue or a different financial-product domain (e.g.\ savings accounts, insurance) with comparatively little modification, since its core method --- extracting named/numeric claims via pattern matching and checking them against a structured ground-truth table --- is domain-agnostic.

\section{Concluding Remarks}

This dissertation demonstrates that a free, open-weight LLM can perform a real, factually-grounded personalised recommendation task at a level worth taking seriously as an alternative to conventional content-based filtering, while also demonstrating, empirically and specifically for this task, why prompting-strategy choice cannot be treated as a minor implementation detail: it is the single largest source of variation in output quality found in this study, and it affects different quality dimensions --- accuracy against a baseline, factual reliability, and perceived explanation quality --- in different, sometimes directly opposing, directions. The clearest illustration of this is also the dissertation's central warning for anyone building or evaluating LLM-based recommenders: the strategy whose explanations read best to an independent judge was, on the same data, the strategy least likely to be right.
